%% =====================================================================
%%  B1-T-06-03 -- what B1 contributes to "The Third-Party Lie Rule"
%%  -------------------------------------------------------------------
%%  LAYER A: APPEND-ONLY. Written once, then left alone. Material found
%%  only here sits in the `unique` slot and is protected by rule R10.
%% =====================================================================
%% @kind: source
%% @id: B1-T-06-03
%% @book: B1
%% @topic: T-06-03
%% @locator: Tactic 1, pp. 15--16
%% @status: distilled
%% @unique: yes
%% @integrated: yes
%% @updated: 2026-09-19

\dossier{B1}{T-06-03}{\bookshort{B1}}{Tactic 1, pp. 15--16}

\begin{distilled}
  \item If you hear a person you trust lie to someone else, cancel the trust.
  \item The reasoning given is that a person who lies to somebody else will lie to you; there is no basis for expecting them to think more of you later when lying suits their purpose.
  \item People live by an internal code; if the code permits lying to an enemy it will permit lying to a friend when that gets them something they want.
  \item Friend and enemy are reclassifiable: a friend who gets in the way becomes an enemy easily, so special status is not durable protection.
  \item Never trust a person who lies to others but says they will not lie to you because you are a friend.
\end{distilled}

\begin{unique}
  \item The internal-code model: honesty is governed by category rather than by person, and categories are cheaply reassigned --- which is why special exemption offers are worthless.
  \item The explicit instruction to withdraw trust on observing a lie directed at a third party, rather than to discount it.
\end{unique}
